\documentclass[10pt,a4paper]{article}
\usepackage[margin=18mm]{geometry}
\usepackage[T1]{fontenc}
\usepackage{lmodern,microtype,amsmath,amssymb,booktabs,tabularx,enumitem}
\usepackage{graphicx,xcolor,caption,fancyhdr}
\usepackage[hidelinks]{hyperref}
\usepackage{xurl}
\definecolor{ink}{HTML}{183344}
\definecolor{teal}{HTML}{237A69}
\definecolor{rust}{HTML}{A34F37}
\definecolor{muted}{HTML}{52616B}
\setlength{\parindent}{0pt}
\setlength{\parskip}{5pt}
\setlist{nosep,leftmargin=*}
\captionsetup{font=small,labelfont=bf,skip=5pt,hypcap=false}
\pagestyle{fancy}\fancyhf{}
\fancyhead[L]{\small\color{muted}FEEC acoustic benchmark / exploratory findings}
\fancyhead[R]{\small\color{muted}8 November 2026}
\fancyfoot[R]{\small\thepage}\renewcommand{\headrulewidth}{0pt}
\setlength{\headheight}{14pt}
\newcommand{\pct}{\%}
\newcommand{\takeaway}[1]{\vspace{3pt}{\large\bfseries\color{ink}#1}\par\vspace{3pt}}
\graphicspath{{feecbench/reports/acoustic-progress/figures/}}
\begin{document}
{\LARGE\bfseries\color{ink}What improves, what breaks down,\\and what remains untested}\par
{\large FEEC acoustic benchmark}\par
\takeaway{Balancing the loss improves velocity prediction.\\Learning the coupled dynamics remains difficult.}
All full-run results below use attention without differential transport. The new
transport extension has implementation checks, but no full-run predictive result.

\renewcommand{\arraystretch}{1.25}
\begin{tabularx}{\linewidth}{@{}>{\raggedright\arraybackslash}p{29mm}X>{\raggedright\arraybackslash}p{34mm}@{}}
\toprule
Question & Evidence & Assessment\\\midrule
Reference solver & Independent backend agreement and convergence; sampled resolution audit & \textcolor{teal}{Supported, with limits}\\
Field weighting & Validation velocity error falls from 92.1\pct\ to 40.4\pct\ with balanced MSE & \textcolor{teal}{Clear improvement}\\
Coupled dynamics & Exact-velocity-family pressure error: 60.9\pct\ train, 67.9\pct\ validation (fixed queries) & \textcolor{rust}{Still limited}\\
Transfer & Higher errors on held-out geometry, genus and meshes; mass loss only & \textcolor{rust}{Degrades; still useful}\\
Typed transport & CPU implementation checks; full experiments pending & \textcolor{muted}{Benefit untested}\\
\bottomrule
\end{tabularx}

\section*{1. What succeeds: correcting field imbalance}
\includegraphics[width=\linewidth]{loss_comparison.pdf}
\captionof{figure}{Regular-query validation, 64 configurations. Dots are seeds 0--2;
horizontal marks and labels are means. All objectives use the same corpus and
25,600 updates. Lower is better; these are errors, not pass rates.}

\textbf{Interpretation.} Ordinary MSE has approximately 173:1 pressure-to-velocity
zero-prediction loss contributions after physical normalization. Balancing these
contributions recovers much of mass weighting's aggregate velocity improvement:
92.11\pct\ (MSE), 39.69\pct\ (mass), 40.36\pct\ (balanced).
This supports field balancing as a required control; it does not establish that
mass geometry is irrelevant or that either objective solves the dynamics.

\textbf{Reading the evidence.} Errors use mass norms in raw physical units,
normalized over each target trajectory, then averaged equally over configurations
and seeds. Fixed-query fitting diagnostics on the next page use a separate query
set. Full-run results are exploratory; numerical scope and methods are in the appendix.

\clearpage
\section*{2. Where it breaks down: family and response field}
\takeaway{Exact-velocity inputs expose a pressure-fitting limitation.}
The family averages separate pressure-driven responses from velocity-driven pressure.
A high training error shows that fitting itself remains difficult; the additional
validation gap measures a different limitation.

\includegraphics[width=\linewidth]{family_response.pdf}
\captionof{figure}{Fixed-query diagnostic at $t=0.26,0.50,0.76$. Each cell is the
three-seed mean error in percent; 128 training and 16 validation configurations per
family. All panels share a 0--110\pct\ scale. Colors encode error magnitude only,
with no universal success threshold. These are whole-field errors within each family,
not isolated causal-response errors.}

\textbf{Read the pressure row for exact velocity.} With balanced MSE, training
pressure error is 60.91\pct\ and validation error is 67.85\pct: a 6.94 percentage-point
gap on top of a large fitting error. The full-run limitation is already present on
seen configurations. It cannot be explained by distribution shift alone.

\textbf{Read velocity across the pulse families.} Balancing markedly improves
velocity compared with ordinary MSE, but pulse-family error remains
36.71\pct\ on training and 37.81\pct\ on validation. Better aggregate accuracy is
not the same as accurate pressure-to-velocity response.

\textbf{Keep stationary and dynamic content separate.} Pulse-plus-harmonic velocity
includes a stationary component. Its lower whole-field error does not by itself
prove better dynamic coupling. In regular-query balanced validation, harmonic
component error in this family averages 12.83\pct. Report exact, harmonic and
spurious coexact content separately; zero-target components require absolute leakage.

\takeaway{This is a limitation of the reported training runs, not proof of architectural impossibility.}
Two-case local fits passed both coupling directions after a longer budget; a
2,000-update cosine run reached 6.30/4.41\pct\ source-response errors. These are
single-seed, grid-12 diagnostics, not directly comparable full-corpus results.

\clearpage
\section*{3. How failure evolves over time}
\takeaway{Field error and conserved quantities answer different questions.}
These two validation examples use balanced MSE. Selection is deterministic: rank
configurations within each family by three-seed mean pressure trajectory error,
choose the nearest-rank median pulse and 90th-percentile exact-velocity case,
and display seed 0. They illustrate behavior rather than establish population rates.

\includegraphics[width=\linewidth]{case_dynamics.pdf}
\captionof{figure}{Positive regular queries in the training horizon $0<t\leq1$.
Top: saved mass-normalized pressure and velocity error curves, with a common vertical
scale. Bottom: predicted and reference total energy relative to reference energy
at the first plotted time. Case IDs and selection ranks are saved in
\texttt{figures/provenance.json}. The typical pulse is a relative success, not a
threshold-qualified pass.}

\textbf{Population context.} Across all balanced validation cases and seeds,
the mean of each configuration's maximum relative energy error is 27.63\pct;
for the exact-velocity family it rises to 66.55\pct. This is a mean of per-case
maxima, not the maximum of a mean energy curve. Lower field error has not
established faithful physics.

\textbf{Read across the cases.} The typical pulse keeps pressure error relatively
low, while velocity error grows. The difficult exact-velocity case has large
pressure error and loses most of its total energy by the end of the horizon.
Energy preservation alone would still be insufficient: frozen
initial fields can preserve total energy while missing the exchange between fields.

\textbf{What remains to visualize.} The available full-run export contains
predictions and scored curves, but not the matching reference meshes/fields or
separate pressure/velocity energy curves. Spatial reference--prediction--error
panels and energy-exchange plots therefore remain pending those inputs. They should
use these same selected cases, three matched times, shared reference/prediction
color scales and an independent error scale. No surrogate spatial data are used here.

\clearpage
\section*{4. How far the improvement extends}
\takeaway{Mass-loss predictions transfer with degradation. Balanced-MSE OOD is untested.}
\includegraphics[width=\linewidth]{transfer.pdf}
\captionof{figure}{Original mass-loss campaign, regular queries, three-seed means,
64 configurations per setting. Labels show absolute error and change from ID medium
in percentage points (pp). Mesh variants are paired to the ID configurations;
geometry/genus shifts compare separate partitions. Dashed lines mark ID medium.}

\textbf{Accuracy context.} Mass-loss ID-medium error is 28.23/38.34\pct\ for
pressure/velocity. The truncated modal control reaches 1.33/1.72\pct\ against the
same-mesh oracle. It uses the known linear operator, is not a learned baseline,
and its score is not continuum accuracy. Neural improvement over persistence and
zero controls does not close this numerical-control gap.

\textbf{Scope.} Genus-3 transfer also changes geometry. These results establish
neither isolated topological generalization nor balanced-MSE OOD performance.
Off-grid validation stays within the same horizon: aggregate pressure/velocity
errors are 30.29/90.37\pct\ (MSE), 28.00/38.82\pct\ (mass), and
27.02/39.46\pct\ (balanced). There is no long-time claim.

\section*{5. The next experiment must distinguish the explanations}
\begin{tabularx}{\linewidth}{@{}>{\raggedright\arraybackslash}p{31mm}XX@{}}
\toprule
Hypothesis & Matched experiment & Evidence needed\\\midrule
Optimization limits fitting & Constant/cosine continuation with matched update budgets & Both coupling responses improve on training and validation\\
Differential transport helps & Attention versus $D_0/\delta_1$ messages; same loss, seeds and queries & Paired family/field gains, physics diagnostics and capacity/compute controls\\
Improvements generalize & Frozen protocol and fresh held-out evaluation & Balanced and transport results across geometry, genus and meshes\\
\bottomrule
\end{tabularx}

\textbf{Numerical prerequisite.} Complete mixed-pressure exact-velocity and remeshing
reference checks before confirmatory claims. The grid-16 audit's largest sampled
velocity difference is 4.99\pct, close to the provisional 5\pct\ margin; backend
agreement does not remove this resolution uncertainty.
\clearpage\section*{Appendix A. Experimental setup and numerical evidence}

\textbf{Task.} Predict pressure $p(t)\in C^0$ and velocity $u(t)\in C^1$ from an
initial state, surface mesh and query time. Connected, closed orientable surfaces
have genus 1--3, constant sound speed $c=1$, no forcing, damping or boundaries,
and horizon $T=1$. The Whitney FEEC reference solves
\[
 M_0\dot p=D_0^\top M_1u,\qquad \dot u=-D_0p,
 \qquad E=\tfrac12(p^\top M_0p+u^\top M_1u).
\]
PyDEC supplies consistent masses; implicit midpoint advances the solution.
The four initial-condition families are one pulse, two pulses, pulse plus harmonic
velocity, and exact velocity $u_0=D_0\psi$ with zero or independently varied pressure.
Harmonic velocity alone is stationary; exact velocity is essential for testing
velocity-driven pressure dynamics.

\textbf{Data and controls.} Train on 512 genus-1/2 configurations; use 64 each for
validation, ID test, geometry OOD and genus-3 OOD. Coarse, fine and independent
remeshes of the ID set give 960 mesh samples from 768 configurations. Train on
medium meshes, grid 16; use $\Delta t=0.002$, 51 regular snapshots and 64 pressure
modes. Positive regular times train the model; off-grid queries are evaluated
separately within the same horizon. Splits group physical configurations and
geometries; scales use training-medium data only. Zero prediction, persistence
and a truncated modal propagator are controls.

\textbf{Model and budget.} The scalar Transolver-style adapter has separate vertex
and edge encoders/heads, shared slice attention, canonical edge signs and
length-aware cochains: four blocks, width 128, 32 slices, 681,602 parameters.
The forward pass uses geometry, initial fields, time and $c$; FEEC operators
originally entered the oracle, losses and diagnostics only. Six mass/MSE runs
and three fresh balanced-MSE runs use paired seeds 0--2, AdamW, constant
learning rate $10^{-3}$, accumulation four, and 200 epochs/25,600 updates each.
Checkpoints share validation mass-score selection. The balanced study reuses the
same corpus and evaluates only training/validation; its recorded GPU is an H100
NVL MIG 3g.47gb. This was intentionally a baseline/loss study, not the complete
Transformer--TNO architecture experiment.


Independent MFEM assembly/stepping passed 36 cases: maximum relative matrix and
trajectory discrepancies were $6.14\times10^{-14}$ and $1.26\times10^{-13}$.
Temporal orders were approximately two; an analytical sphere test gave spatial/
geometric orders 1.82--1.89. These validate different aspects of the oracle;
backend agreement and energy conservation alone do not bound discretization error.

The older grid-12 development study failed its provisional middle-band gate
(6.51\pct) and reached 9.87\pct\ velocity discrepancy. The newer disjoint,
24-configuration grid-16 audit has maximum spatial-plus-geometry differences
of 0.94\pct\ low-band pressure, 4.18\pct\ middle-band pressure and 4.99\pct\ velocity.
The largest temporal velocity difference is 0.52\pct. These sampled fine-grid
differences support exploratory comparisons but are not continuum bounds or a
frozen protocol; the audit exact-velocity cases have zero initial pressure.

Local two-case grid-12 diagnostics also caution against declaring an architecture
failure: 500-update fits failed the 10\pct\ coupling criterion, whereas a balanced
scalar fit at 1,500 updates passed both directions; 2,000 updates with cosine
decay reached 6.30/4.41\pct\ source-response errors. These are single-seed fitting
results. Changing both schedule and budget does not isolate either effect.
Directional-feature additions did not establish a general fix. Matched full-run
constant/cosine continuation is prepared, but no full-run result is included here.

\clearpage\section*{Appendix B. Typed transport implementation}

The optional \texttt{--transport d0\_delta1} adds simultaneous residual differential
messages before every attention block. With canonical edge-density hidden features
$Z_1$, requested-orientation signs $S$ and edge lengths $L$, the messages are
\[
 m_{0\to1}=L^{-1}S D_0Z_0,\qquad
 m_{1\to0}=M_0^{-1}D_0^\top M_1SLZ_1.
\]
The $SL$ factors convert densities to oriented integral cochains and back.
Hidden channels are learned features rather than physical $p/u$; physical field
scale ratios are not inserted into these latent maps. Each message receives
channel LayerNorm and a learned projection before the residual update.

The full consistent $M_0$ system uses device-local, multi-column preconditioned
CG and an implicit adjoint solve in backward, with no dense inverse. Jacobi is
only the preconditioner. Float32/float64 relative tolerances are $10^{-6}/10^{-10}$,
with at most 64 iterations and a true-residual check allowing five times that
tolerance; failed solves raise. Mesh/operator differentiation and mixed precision
are unsupported. Output projections start at zero, preserving shared baseline
initialization and initial predictions. The standard variant has 814,722 parameters
(+133,120, approximately 19.5\pct); equal updates do not equalize capacity or compute.

Transport is recorded in training/checkpoint identities and restored for evaluation
and diagnostics. Resume rejects architecture changes; the default
\texttt{--transport none} preserves the baseline. Historical campaign source guards
remain active, so new source/architecture runs require fresh output directories.
Loss-only comparison tools reject mixed transports. CPU checks cover operator
agreement, adjoint gradients, active orientation flips, initial-state recovery,
both coupling-direction overfits, and training/resume/evaluation integration.
The full CPU suite passes 111 tests, with three CUDA/backend checks skipped;
Ruff and the three-epoch command-line CPU smoke workflow also pass.
These are implementation checks, not predictive-performance results.
\textbf{No full GPU transport or OOD result is claimed.}

\section*{Interpretation and remaining limitations}

\begin{enumerate}
\item Keep balanced MSE as a required control; separate field balancing, geometric
loss weighting and architectural transport. Next compare attention with transport
using the same corpus, objective, paired seeds, query sampling, selection score and
budget. Add capacity/compute controls before attributing gains solely to transport.
\item Report pressure/velocity by family, genus and mesh, plus exact/harmonic/coexact
errors, energy, residuals and time curves. Zero-target harmonic content needs absolute
leakage scores, not floor-dominated relative errors. Residual evaluation currently
covers one sample per run, so it cannot establish corpus-wide PDE fidelity.
\item Complete reference checks for mixed-pressure exact-velocity cases and remeshing;
review propagation horizon, spectral eligibility and numerical margins before
protocol freeze or spectral-loss claims. The 5\pct\ velocity margin is very narrow.
\item Genus transfer is confounded with geometry. All four IC families occur in
training; no unseen-family, coefficient, long-time, autoregressive or cross-PDE
transfer is tested. Harmonic ICs span only a $g$-dimensional subspace of the full
$2g$ space; reference velocity has no coexact component. General commuting transfers
between curved/remeshed spaces are not established.
\item This is an acoustic $C^0/C^1$ transport extension, not the full Transformer--TNO
foundation model. Canonicalization handles edge flips, but the nonlinear architecture
does not guarantee conservation or intrinsic equivariance. CUDA transport agreement,
runtime, memory and predictive benefit remain unmeasured. Existing held-out scores
informed these design decisions; a confirmatory claim needs a frozen protocol and
fresh evaluation evidence. The modal control remains essential for accuracy/cost context.
\end{enumerate}

\clearpage\section*{Appendix C. Evidence and reproduction}

\small
Repository-relative sources: \path{plan.md}; \path{feecbench/SURFACE_EXPERIMENT.md};
\path{feecbench/DECISIONS.md}; full campaign export
\path{reports/acoustic/analysis.json}; latest study
\path{reports/acoustic/loss-balance-g16/analysis.json} and its per-seed
\path{runs/*/evaluation/per_sample.json}; independent validation
\path{feecbench/reports/surface_validation_mfem.json}.
Original proposals: \textit{PDE Benchmark}, Sections 1.3, 4--6, and
\textit{Typed Physics Benchmark and Transformer--TNO}, Sections 8--15.
Full-run results date from 5--7 November; transport implementation is dated
8 November. Generated corpora and reports remain outside version control.

From \path{feecbench/}, a small end-to-end check is
\texttt{bash scripts/run\_acoustic\_baseline.sh smoke --device cpu}
\texttt{--transport d0\_delta1 --objective mse\_balanced}.
Training/diagnostic commands for the existing full corpus are documented in
\path{SURFACE_EXPERIMENT.md}. To regenerate the figures and compile this source from the repository root:
\begin{verbatim}
python3 feecbench/scripts/plot_acoustic_progress.py
pdflatex -output-directory=feecbench/reports/acoustic-progress \
  feecbench/docs/acoustic_progress.tex
\end{verbatim}
Run LaTeX twice to resolve table references.

\normalsize
\textbf{Figure provenance.} The plotting script reads the original campaign summary,
the matched loss study and all three balanced per-sample evaluation files. It checks
campaign identity, objective/seed coverage and aggregation counts. It writes source
SHA-256 hashes, unrounded plotted seed means and deterministic case-selection details
to \path{feecbench/reports/acoustic-progress/figures/provenance.json}.
Figures and compiled outputs are generated under the ignored reports directory.
The report does not rerun training or reference simulations.

\textbf{Comparison detail.} Regular-query validation mean $\pm$ seed standard deviation
(not a confidence interval): MSE pressure $31.11\pm2.76$\pct, velocity
$92.11\pm1.85$\pct; mass pressure $28.61\pm2.53$\pct, velocity $39.69\pm2.87$\pct;
balanced pressure $27.60\pm1.52$\pct, velocity $40.36\pm4.40$\pct.
Training-derived balancing weights are 0.502886 for pressure and 87.1354 for velocity;
neither floor activates. The common multiplier matches expected zero-prediction
loss, not gradients, clipping or optimizer trajectories. Balanced MSE improves
fixed-query validation velocity over ordinary MSE in every family/seed, but does
not uniformly beat mass loss.
\end{document}
